\documentclass[aspectratio=169,11pt,usenames,dvipsnames]{beamer}

\usepackage[czech,english]{babel}
\usepackage{CJKutf8}
\newcommand{\zh}[1]{\begin{CJK}{UTF8}{gbsn}#1\end{CJK}}
\usepackage{graphicx}
\usepackage{enumitem}
\usepackage{amsmath}
\newlength{\dhatheight}
\newcommand{\doublehat}[1]{%
    \settoheight{\dhatheight}{\ensuremath{\hat{#1}}}%
    \addtolength{\dhatheight}{-0.2ex}%
    \hat{\vphantom{\rule{1pt}{\dhatheight}}%
    \smash{\hat{#1}}}}
\usepackage{mathtools}
\usepackage{float}
\usepackage{tikz}
\usetikzlibrary{patterns,arrows.meta,calc,3d,angles}
\usepackage{tkz-euclide}
\tikzset{point style/.style = {%
  draw = black,
  inner sep = 0pt,
  shape = circle,
  minimum size = 5pt,
  fill = black
 }
}
\usepackage{enumitem}

\usepackage{caption}
\usepackage{subcaption}

\usepackage{booktabs}
% Flowchart stuff

\usepackage{pgfopts}
\usepackage{xcolor}
\usepackage{tcolorbox}

\usetheme[
 titlestyle=style2,
 titleformat=smallcaps,
 sectionstyle=plain,
 slidestyle=cyber,
 headingcolor=theme,
 block=transparent
]{trigon}

\title{Prime Factorizations\\
GCD \& LCM}
\date{\today}
\author{Adam Klepáč}
\institute[GEVO]{Gymnázium Evolution Jižní Město}
\biglogo[width=.2\textwidth]{logo}
\smalllogo[width=.1\textwidth]{logo}
\titlegraphic{\includegraphics[height=\paperheight]{title.jpg}}

\def\subsectionname{}

% enumerate and itemize global settings
\setlist{topsep=0pt}
\setlist[itemize,1]{label=\textbullet}
\setlist[itemize,2]{label=\textopenbullet}
\setlist[enumerate,1]{label=\arabic*.}
\setlist[enumerate,2]{label=\alph*)}

% custom colors %
\definecolor{Licorice}{HTML}{200F08}
\definecolor{Bistre}{HTML}{3D1F0F}
\definecolor{SealBrown}{HTML}{561F0D}
\definecolor{EerieBlack}{HTML}{1A1A17}
\definecolor{EarthYellow}{HTML}{D7AC6C}

\colorlet{tPrim}{Licorice}
\colorlet{tTheme}{SealBrown}
\colorlet{tSec}{EarthYellow}
\colorlet{tAccent}{Bistre}
\colorlet{tText}{EerieBlack}

\newcommand{\clr}{\textcolor{BrickRed}}
\newcommand{\clb}{\textcolor{RoyalBlue}}
\newcommand{\clg}{\textcolor{ForestGreen}}
\newcommand{\clm}{\textcolor{Magenta}}
\newcommand{\cls}{\textcolor{Salmon}}
\newcommand{\clo}{\textcolor{OrangeRed}}

\newcommand{\N}{\mathbb{N}}
\newcommand{\Z}{\mathbb{Z}}
\renewcommand{\P}{\mathbb{P}}
\DeclareMathOperator{\tng}{\triangle}
\DeclareMathOperator{\dv}{div}

\tcbset{
 boxsep=7pt,
 fonttitle=\sc,
 colframe=tGreyBg,
 colframe=tSec,
 boxrule=1pt
}

\begin{document}
\titleframe

\begin{frame}
 \frametitle{Contents}
 \tableofcontents
\end{frame}

\section{Divisibility}

\begin{frame}
 \frametitle{Divisibility}
  We say that a number $n \in \N$ \alert{is divisible} by $m \in \N$ or that $m$
  \alert{divides} $n$ \pause
  if there exists a number $q \in \N$ (called the \alert{quotient}) such that
  \[
   n = m \cdot q.
  \]
  \pause
  We write the fact that $m$ divides $n$ as $m \mid n$.
\end{frame}

\begin{frame}
 \frametitle{Divisibility}
 Even if $m$ doesn't divide $n$, we can always at least \alert{divide with
 remainder}.\\
 \pause
 Formally, for every two numbers $n,m \in \N$, there always exist number $q \in
 \N$ and $r \leq n$ (called the \alert{remainder}) such that
 \[
  n = m \cdot q + r.
 \]
\end{frame}

\begin{frame}
 \frametitle{Integer Division \& Modulus}
 \begin{itemize}
  \item The operation which to a pair $(n,m)$ assigns the \alert{quotient} of
   the division of $n$ by $m$ is called \alert{integer division} and denoted by
   $\dv$.
  \pause
  \begin{itemize}
   \item For example, $7 \dv 2 = 3$ and $5 \dv 9 = 0$.
  \end{itemize}
  \pause
 \item The operation which to a pair $(n,m)$ assigns the \alert{remainder} of
  the division of $n$ by $m$ is called \alert{modulus} and denoted $\bmod$.
  \pause
  \begin{itemize}
   \item For example, $7 \bmod 2 = 1$ and $5 \bmod 9 = 5$.
  \end{itemize}
 \end{itemize}
\end{frame}

\begin{frame}
 \frametitle{Integer Division \& Modulus}
 Using integer division and modulus, we can write the operation of division with
 remainder like this:
 \[
  n = m \cdot (n \dv m) + n \bmod m.
 \]
\end{frame}

\section{Chinese Remainder Theorem}

\begin{frame}
 \frametitle{Counting Soldiers}
 The famous Chinese military general, strategist and philosopher \zh{孙子}
 describes in his book \zh{孙子兵法} an efficient method to count the number of
 soldiers.\\
 \pause
 He instructs the soldiers to arrange into a grid with
 \begin{itemize}
  \item 21 soldiers in each row. He counts the number of soldiers remaining in
   the last row to be 3.
  \pause
  \item 25 soldiers in each row. He counts the number of soldiers remaining in
   the last row to be 6.
  \pause
  \item 32 soldiers in each row. He counts the number of soldiers remaining in
   the last row to be 23.
 \end{itemize}
 \pause
 \zh{孙子} then determines, that there are 1431 soldiers standing before him.
\end{frame}

\subsection{Prime Factorization}

\begin{frame}
 \subsectionpage
\end{frame}

\begin{frame}
 \frametitle{Prime Number}
 \begin{tcolorbox}[title=Prime Number]
  A number $n \in \N$ is \alert{prime} if it has \alert{exactly two divisors} --
  $1$ and itself.

  We denote the set of all prime numbers by $\mathbb{P}$.
 \end{tcolorbox}
\end{frame}

\begin{frame}
 \frametitle{Decomposition Into Primes}
 Prime numbers are somewhat of a building blocks of other numbers.\\
 \pause
 Each natural number decomposes \alert{uniquely} into a product of prime
 numbers.\\
 \pause
 Formally,
 \begin{tcolorbox}[title=Prime Factorization]
  For each number $n \in \N$, there exist prime numbers $p_1,\ldots,p_m \in \P$ and
  powers $k_1,\ldots k_m \in \N$ such that
  \[
   n = p_1^{k_1} \cdot p_2^{k_2} \cdot \ldots \cdot p_m^{k_m}.
  \]
 \end{tcolorbox}
 \pause
 The numbers $p_1,\ldots,p_m$ are called the \alert{prime factors} of $n$.
\end{frame}

\begin{frame}
 \frametitle{Decomposition Into Primes}
 Currently, there's no known fast algorithm how to decompose a number into its
 prime factors.\\
 \pause
 The best we can do is to try each prime number from $2$ up to $\sqrt{n}$.\\
 \pause
 We divide the number in question as many times as possible and proceed to find
 the next prime factor of the resulting quotient.
\end{frame}

\begin{frame}
 \frametitle{Decomposition Into Primes -- Example}
 Let's decompose the number 16335 into its prime factors.\\
 \pause
 It's as easy as keeping a table of all the prime factors we've divided the
 number by.\\
 \pause
 \begin{center}
  \begin{tabular}{c|c}
   \textbf{Prime Factors} & \textbf{Quotients}\\
   \toprule
   3 & 5445\\
   \pause
   3 & 1815\\
   \pause
   3 & 605\\
   \pause
   5 & 121\\
   \pause
   11 & 11\\
   \pause
   11 & 1
  \end{tabular}
 \end{center}
 \pause
 This means that $16335 = 3^3 \cdot 5 \cdot 11^2$.
\end{frame}

\begin{frame}
 \frametitle{Greatest Common Divisor}
 \begin{tcolorbox}[title=Greatest Common Divisor]
  Given two numbers $n,m \in \N$, their \alert{Greatest Common Divisor} (GCD) or
  Highest Common Factor (HCF) is the largest natural number $k \in \N$ such that
  $k \mid n$ and $k \mid m$.
 \end{tcolorbox}
\end{frame}

\begin{frame}
 \frametitle{Greatest Common Divisor}
 The simplest (and most inefficient) way how to compute the GCD is to decompose
 the two numbers into prime factors and take all those that the two numbers have
 in common.\\
 \pause
 For example, let's take a look at the decompositions of 16335 and 17325.\\
 \vspace{1em}
 \centering
 \begin{minipage}{.48\textwidth}
  \centering
   \begin{tabular}{c|c}
    \textbf{Prime Factors} & \textbf{Quotients}\\
    \toprule
    3 & 5445\\
    3 & 1815\\
    3 & 605\\
    5 & 121\\
    11 & 11\\
    11 & 1
  \end{tabular}
 \end{minipage}
 \begin{minipage}{.48\textwidth}
   \begin{tabular}{c|c}
    \textbf{Prime Factors} & \textbf{Quotients}\\
    \toprule
    3 & 5775\\
    3 & 1925\\
    5 & 385\\
    5 & 77\\
    7 & 11\\
    11 & 1
  \end{tabular}
 \end{minipage}
\end{frame}

\begin{frame}
 \frametitle{Greatest Common Divisor}
 If we highlight the prime factors the two numbers share\\
 \vspace{1em}
 \begin{minipage}{.48\textwidth}
  \centering
   \begin{tabular}{c|c}
    \textbf{Prime Factors} & \textbf{Quotients}\\
    \toprule
    \clr{3} & 5445\\
    \clr{3} & 1815\\
    3 & 605\\
    \clg{5} & 121\\
    \clb{11} & 11\\
    11 & 1
  \end{tabular}
 \end{minipage}
 \begin{minipage}{.48\textwidth}
   \begin{tabular}{c|c}
    \textbf{Prime Factors} & \textbf{Quotients}\\
    \toprule
    \clr{3} & 5775\\
    \clr{3} & 1925\\
    \clg{5} & 385\\
    5 & 77\\
    7 & 11\\
    \clb{11} & 1
  \end{tabular}
 \end{minipage}\\
 \pause
 \vspace{1em}
 we deduce that the GCD of 16335 and 17325 is
 \[
  3^2 \cdot 5 \cdot 11 = 495.
 \]
 \pause
 We'll denote the GCD of $n$ and $m$ as $\mathrm{gcd}(n,m)$ or simply as
 $(n,m)$.
\end{frame}

\begin{frame}
 \frametitle{Least Common Multiple}
 \begin{tcolorbox}[title=Least Common Multiple]
  Given $n,m \in \N$ their \alert{Least Common Multiple} (LCM) is the
  \alert{smallest} number $k \in \N$ such that $n \mid k$ and $m \mid k$.
  \pause
  It is given by the formula
  \[
   \mathrm{lcm}(n,m) = \frac{n \cdot m}{\mathrm{gcd}(n,m)}.
  \]
 \end{tcolorbox}
\end{frame}

\begin{frame}
 \frametitle{Least Common Multiple}
 The other option how to calculate the least common multiple of two numbers is
 to take their prime factors elevated on the highest power that appears in at
 least one of the numbers.
\end{frame}

\begin{frame}
 \frametitle{Least Common Multiple}
 Explicitly, from the prime decompositions of 16335 and 17325\\
 \vspace{1em}
 \begin{minipage}{.48\textwidth}
  \centering
   \begin{tabular}{c|c}
    \textbf{Prime Factors} & \textbf{Quotients}\\
    \toprule
    \clr{3} & 5445\\
    \clr{3} & 1815\\
    \clr{3} & 605\\
    5 & 121\\
    \clb{11} & 11\\
    \clb{11} & 1
  \end{tabular}
 \end{minipage}
 \begin{minipage}{.48\textwidth}
   \begin{tabular}{c|c}
    \textbf{Prime Factors} & \textbf{Quotients}\\
    \toprule
    3 & 5775\\
    3 & 1925\\
    \clg{5} & 385\\
    \clg{5} & 77\\
    \clm{7} & 11\\
    11 & 1
  \end{tabular}
 \end{minipage}\\
 \pause
 \vspace{1em}
 we see that their LCM is
 \[
  3^3 \cdot 5^2 \cdot 7 \cdot 11^2 = 571725.
 \]
\end{frame}

\end{document}
